\section{Mechanism-Gated Self-Evolution}
\label{sec:method}

\subsection{Immutable candidate lifecycle}

Each round executes four stages: observe, diagnose, mutate, and verify.
Exploration trajectories are immutable and role-filtered.  A diagnosis cites
specific trajectory spans, localizes a failure to an editable surface, and
predicts both an intended effect and plausible regressions.  A typed mutation
is applied to exactly one parent snapshot, structurally validated, and assigned
a content hash before execution.  The proposer cannot rewrite an earlier
artifact or its own evaluation policy.

The same frozen model configuration may fill solver and self-evolution roles;
all role calls are recorded separately.  A claim of \emph{self-evolution}
requires that this model family closes execution, diagnosis, proposal, and
selection without manual candidate choice.  An external stronger proposer is a
valid harness-optimization baseline, but answers a different question.

\subsection{Matched interventions}

For each gate nomination, the verifier constructs a matched quartet:
\begin{enumerate}
  \item \textbf{child}: the nominated harness with the mutation enabled;
  \item \textbf{parent}: the exact champion from which the child descends;
  \item \textbf{revert}: the child runtime with the mutation exactly disabled,
  retaining other candidate-path overhead; and
  \item \textbf{placebo}: a format- and budget-matched neutral mutation that
  should not encode the proposed mechanism.
\end{enumerate}
Task identity, rollout seed, model settings, runtime image, and resource ceiling
are paired.  Arm order is randomized by a verifier-owned schedule.  Revert
tests whether the declared edit matters; placebo tests whether merely adding a
similarly shaped artifact or execution path explains the observation.

\subsection{Fail-closed promotion gate}

The gate evaluates source-linked claims in causal order.  Later success cannot
repair an earlier missing link:
\[
\promote(c)=
\mathbb{1}\{A_c\land D_c\land B_c\land U_c\land P_c\land C_c\land L_c\},
\]
where \(A,D,B\) denote activation, adherence, and intended behavior; \(U\)
denotes utility benefit; \(P\) protected-slice non-inferiority; \(C\) the cost
constraint; and \(L\) protocol, provenance, and leakage validity.  Every term is
three-valued---pass, fail, or inconclusive---and only all-pass promotes.  Safety
and cost remain separate constraints rather than disappearing inside a weighted
score.

The verifier spends a frozen campaign-level error budget across all possible
nominations.  A conservative fixed-horizon allocation is the default; a
time-uniform procedure may replace it only when preregistered before outcomes
are visible \citep{howard2021confidence}.  The optimizer sees only evidence
authorized for its condition and never sealed outcomes.  Rejected and
inconclusive candidates remain in the ledger, so failures cannot vanish through
complete-case reporting.

\subsection{Independent selection and sealing}

The promotion authority is outside every candidate's editable surface.  It
owns task manifests, paired schedules, graders, evidence producers, usage
accounting, and the champion pointer.  After every condition and independent
search run reaches a terminal selection state, a separate sealed service accepts
only authorized champion hashes.  One atomic authorization evaluates the full
grid and releases signed aggregates; it never returns item outputs or
trajectories to the optimizer.  Opening any sealed item early invalidates that
study lineage rather than turning the item into more validation data.
